\section{Positive-time expansion at \(u_{205}\)}
\label{sec:positive-time-spike}
The second-order disk of Section~\ref{sec:orbit-orthogonal-kick-spike}
is silent at \(n=205\). A positive-time diagram expansion that keeps
trace orthogonality still clears the Abel threshold at that one dimension.
The canonical proof is
\path{mathematics/fixed-map/positive_time_spike.md}.
The integer comparisons are certified by
\path{formalization/Bgs/PositiveTimeRational.lean} and recomputed by the
Python and C++ verifiers named there. No external refereeing is asserted.
BGS-TRACE-005 and BGS-BGS-001 remain open.

\begin{theorem}[Positive-time spike, BGS-TRACE-119]
\label{thm:positive-time-spike}
Let \(n=205\) and \(N=u_n\), with \(u_m\) the Pell companion of
Section~\ref{sec:kick-blind-arithmetic-spike}. At the fixed coupling
\(\tau=10^{-80}\),
\[
 L_N(U_N;1)>\frac34.
\]
The distance from this value to \((1-e^{-2})/4\) is greater than \(1/2\).
The integer \(u_{205}\) has \(79\) decimal digits, so any statement that the
distance is less than \(1/2\) for every odd \(N\ge N_0\) requires
\(N_0>u_{205}\). At \(n=207\) the pure second-order center satisfies
\((1-\varepsilon_2)^2<3/4\), so no disk of radius \(E\ge0\) about
\(N(1-\varepsilon_2)\) proves the gap there. The estimate does not decide
the limit or a generic odd dimension.
\end{theorem}
